% TEMPLATE-MILC-v3.tex - plantilla maestra UNICA de las guias MILC v3.
%
% La usan las tres familias de guias, cada una con su driver:
%   · Grados 6-11  -> scripts/build-guias-g11.py
%   · Bebras       -> scripts/build-guias-bebras.py
%   · Semillero    -> scripts/build-guias-semillero.py
%
% Antes habia tres copias casi identicas (~400 lineas iguales, 6 distintas):
% cada mejora habia que aplicarla tres veces, y en la practica no se hacia
% (el motor de recursos vivia solo en dos de ellas). Lo que varia entre
% familias esta parametrizado en 6 placeholders que cada driver llena
% (se nombran aqui SIN su sintaxis real: el builder reemplaza por texto plano,
% tambien dentro de los comentarios, y un valor multilinea rompe el preambulo):
%
%   HEADER_IZQ        encabezado izquierdo de pagina
%   HEADER_DER        encabezado derecho de pagina
%   PORTADA_ETIQUETA  rotulo sobre el numero grande (GRADO/MOMENTO/MODULO)
%   PORTADA_SERIE     linea de serie de la portada
%   PORTADA_ID        identificador de la guia en la portada
%   PORTADA_CIRCULO   el circulito de grado (°), o vacio si no aplica
%
% El resto de claves las documenta content/guias/_SCHEMA.md. El script valida
% que no queden placeholders sin reemplazar antes de escribir el archivo final.

\documentclass[11pt,letterpaper]{article}
\usepackage[margin=1.75cm]{geometry}
\usepackage[table]{xcolor}
\usepackage{fontspec}
\usepackage[spanish]{babel}
\usepackage{tikz}
% Librerías TikZ para los diagramas de las guías (bloque `recursos:`):
%  · babel  → babel en español vuelve activos caracteres como «>», y TikZ los usa
%             en sus opciones (>=stealth, ->). Sin esta librería, cualquier
%             diagrama con flechas revienta con «\language@active@arg> has an
%             extra }». Debe ir ANTES de que se dibuje nada.
%  · positioning        → sintaxis «below left=2pt and 3pt».
%  · decorations.pathreplacing → llaves ({brace}) para señalar rangos.
\usetikzlibrary{babel,positioning,decorations.pathreplacing}
\usepackage{graphicx}
% Circuitos: el trabajo de electrónica se hace hoy con micro:bit y sus sensores
% integrados (MakeCode, sin cableado), así que no se necesitan esquemáticos.
% Si algún día entran montajes con protoboard, basta descomentar esta línea:
% los diagramas .tex/.tikz declarados en `recursos:` ya se incrustan solos.
% \usepackage{circuitikz}
\usepackage[most]{tcolorbox}
\usepackage{tabularx}
\usepackage{array}
\usepackage{fancyhdr}
\usepackage{titlesec}
\usepackage{ragged2e}
\usepackage{enumitem}
\usepackage{microtype}

% Helvetica Neue no trae la flecha U+2192, asi que cada «→» escrito literalmente
% en el YAML se compilaba a NADA, en silencio (462 flechas en 61 guias: rutas del
% tipo «paleta → tipografia → lockup» salian rotas). Mapeamos el caracter a su
% version matematica, que si tiene glifo. Asi el autor puede seguir escribiendo
% «→» comodamente en el YAML.
\usepackage{newunicodechar}
\newunicodechar{→}{\ensuremath{\rightarrow}}

\setmainfont{Helvetica Neue}
\setsansfont{Helvetica Neue}
\setlength{\parindent}{0pt}
\setlength{\parskip}{6pt}
\hyphenpenalty=200
\exhyphenpenalty=200
\pretolerance=400
\tolerance=2000
\setlength{\emergencystretch}{1em}
\renewcommand{\arraystretch}{1.25}
\setlist{leftmargin=5mm,itemsep=1mm,topsep=1mm}
\newcolumntype{Y}{>{\RaggedRight\arraybackslash}X}

\definecolor{milcMagenta}{HTML}{E6007E}
\definecolor{milcVino}{HTML}{5A0038}
\definecolor{milcTurquesa}{HTML}{0093A5}
\definecolor{milcVerde}{HTML}{58A923}
\definecolor{milcMostaza}{HTML}{E5B400}
\definecolor{milcOcre}{HTML}{B86600}
\definecolor{milcNegro}{HTML}{241816}
\definecolor{milcGris}{HTML}{F7F7F4}
\definecolor{milcLinea}{HTML}{E8E2DD}
\definecolor{milcGradePrimary}{HTML}{0066FF}
\definecolor{milcGradeDark}{HTML}{003D99}
\definecolor{milcGradeSoft}{HTML}{D6E8FF}
\definecolor{milcGradeAccent}{HTML}{CFE7FF}
\definecolor{milcGradeText}{HTML}{FFFFFF}
\color{milcNegro}

\pagestyle{fancy}
\fancyhf{}
\lhead{\small Tecnología e Informática · Grado 11}
\rhead{\small Guía 22 · MILC v3}
\cfoot{\small\thepage}
\renewcommand{\headrulewidth}{0pt}

\newtcolorbox{titlebox}[1]{
  enhanced,colback=#1,colframe=#1,arc=7mm,boxrule=0pt,
  left=7mm,right=7mm,top=3mm,bottom=3mm,
  before skip=8pt,after skip=8pt,
  fontupper=\sffamily\bfseries\Large\color{white}
}

\newtcolorbox{softbox}[3]{
  enhanced,colback=#2,colframe=#1,borderline west={4pt}{0pt}{#1},
  arc=2mm,boxrule=.4pt,left=4mm,right=4mm,top=3mm,bottom=3mm,
  before skip=6pt,after skip=8pt,title={#3},coltitle=white,
  colbacktitle=#1,fonttitle=\sffamily\bfseries,fontupper=\RaggedRight,
  attach boxed title to top left={xshift=3mm,yshift=-2mm},
  boxed title style={arc=2mm, boxrule=0pt}
}

\newtcolorbox{drawbox}[2]{
  enhanced,colback=white,colframe=#1,arc=4mm,boxrule=1pt,
  left=5mm,right=5mm,top=4mm,bottom=4mm,before skip=8pt,after skip=8pt,
  title={#2},coltitle=white,colbacktitle=#1,fonttitle=\sffamily\bfseries,
  fontupper=\RaggedRight,
  attach boxed title to top left={xshift=4mm,yshift=-2mm},
  boxed title style={arc=3mm, boxrule=0pt}
}

\newtcolorbox{infoband}[1]{
  enhanced,colback=#1!8,colframe=#1!50,arc=2mm,boxrule=.5pt,
  left=4mm,right=4mm,top=2mm,bottom=2mm,before skip=4pt,after skip=4pt,
  fontupper=\small\RaggedRight
}

% Puente narrativo entre secciones (contrato editorial Conéctate).
% Da continuidad: "Acabas de X. Ahora vas a Y porque Z."
% Visualmente: banda izquierda en color del grado, texto en itálica pequeña.
\newtcolorbox{puentebox}{
  enhanced,colback=milcGradeSoft,colframe=milcGradeDark,
  borderline west={3pt}{0pt}{milcGradeDark},
  arc=0mm,boxrule=0pt,left=5mm,right=4mm,top=2.5mm,bottom=2.5mm,
  before skip=4pt,after skip=8pt,
  fontupper=\itshape\small\color{milcNegro}\RaggedRight
}
% La flecha va en modo matematico a proposito: Helvetica Neue NO tiene el glifo
% U+2192, asi que \textrightarrow se compilaba a NADA («Missing character: There
% is no → in font Helvetica Neue») y el puente salia sin flecha. En modo
% matematico la flecha viene de la fuente matematica y si se dibuja.
\newcommand{\puente}[1]{%
  \begin{puentebox}%
    \textcolor{milcGradeDark}{$\rightarrow$}\hspace{2mm}#1%
  \end{puentebox}%
}

\newcommand{\linead}[1]{\noindent\color{milcLinea}\rule{#1}{0.5pt}\color{milcNegro}}
\newcommand{\respuesta}[1]{\par\vspace{2mm}\foreach \n in {1,...,#1}{\noindent\color{milcLinea}\rule{\linewidth}{0.5pt}\par\vspace{5mm}}\color{milcNegro}}
\newcommand{\checkbox}{\textcolor{milcGradeDark}{\fbox{\phantom{x}}}\hspace{5pt}}

\newcommand{\phasecard}[4]{
\begin{tcolorbox}[enhanced,colback=#3,colframe=#2,arc=3mm,boxrule=.5pt,height=3.05cm,valign=center,halign=center,left=1.5mm,right=1.5mm,top=2mm,bottom=2mm]
{\sffamily\bfseries\fontsize{22}{24}\selectfont\textcolor{#2}{#1}}\par
{\sffamily\bfseries\small #4}
\end{tcolorbox}
}

% ── Recursos visuales de la guía ──────────────────────────────────────
% Declarados en el bloque `recursos:` del YAML y emitidos por el builder.
% \guiaFigura   → imagen rasterizada o PDF (foto, carta celeste, montaje).
% \guiaDiagrama → fuente TikZ/circuitikz (.tex) incrustada como vector:
%                 es la vía para circuitos y esquemáticos editables.
% #1 = ruta relativa al .tex · #2 = pie (puede ir vacío)
\newcommand{\guiaFigura}[2]{%
\begin{center}
\includegraphics[width=0.88\linewidth,height=0.38\textheight,keepaspectratio]{#1}\\[1.5mm]
{\footnotesize\itshape\textcolor{milcNegro!75}{#2}}
\end{center}
}

\newcommand{\guiaDiagrama}[2]{%
\begin{center}
\input{#1}\\[1.5mm]
{\footnotesize\itshape\textcolor{milcNegro!75}{#2}}
\end{center}
}

\begin{document}
\thispagestyle{empty}
\begin{tikzpicture}[remember picture,overlay]
  \fill[white] (current page.south west) rectangle (current page.north east);
  \path[fill=milcGradePrimary,rounded corners=16mm]
    ([xshift=.55cm,yshift=-.55cm]current page.north west) rectangle
    ([xshift=-.55cm,yshift=3.65cm]current page.south east);

  \node[anchor=north west,text=milcGradeText] at ([xshift=2.0cm,yshift=-1.85cm]current page.north west)
    {\sffamily\bfseries\fontsize{14}{16}\selectfont GRADO};
  \node[anchor=north west,text=milcGradeText,opacity=.82] at ([xshift=2.0cm,yshift=-2.75cm]current page.north west)
    {\sffamily\bfseries\fontsize{9}{11}\selectfont SERIE GUÍAS · TECNOLOGÍA E INFORMÁTICA};

  \begin{scope}[shift={([xshift=-3.05cm,yshift=-2.55cm]current page.north east)}]
    \fill[white,opacity=.80] (-.62,-.52) rectangle (.62,.52);
    \draw[milcGradeText,line width=1pt] (-.62,-.52) rectangle (.62,.52);
    \fill[milcGradeDark] (-.42,-.38) rectangle (-.22,.06);
    \fill[milcGradeAccent] (-.10,-.38) rectangle (.10,.24);
    \fill[milcGradePrimary!60!black] (.22,-.38) rectangle (.42,.38);
  \end{scope}

  \node[anchor=west,text=milcGradeText] at ([xshift=1.9cm,yshift=-12.85cm]current page.north west)
    {\sffamily\bfseries\fontsize{124}{124}\selectfont 11};
  % El circulito de grado (°) solo aplica a las guias de grado («11°»); en
  % Bebras y Semillero el numero grande es un momento/modulo, no un grado.
  \draw[milcGradeText,line width=1.7pt]
    ([xshift=4.52cm,yshift=-11.68cm]current page.north west) circle (.13cm);

  \node[anchor=west,text=milcGradeText,text width=8.7cm,align=left] at ([xshift=10.35cm,yshift=-11.6cm]current page.north west)
    {
     {\sffamily\bfseries\fontsize{19}{22}\selectfont Validación ---\\hablar antes\\de construir}\\[4mm]
     {\sffamily\bfseries\fontsize{23}{26}\selectfont Guía 22-11-TIC}\\[2mm]
     {\sffamily\fontsize{13}{16}\selectfont Período 3 · Proyecto final emprendedor}\\[5mm]
     {\sffamily\bfseries\fontsize{11}{13}\selectfont Institución Educativa Sor María Juliana}\\[1mm]
     {\sffamily\fontsize{10}{12}\selectfont Autor: Dr. Álvaro Cárdenas Orozco}
    };

  \path[fill=milcGradeSoft,rounded corners=7mm]
    ([xshift=1.35cm,yshift=.85cm]current page.south west) rectangle
    ([xshift=-1.35cm,yshift=4.25cm]current page.south east);
  \draw[milcGradeDark,line width=.65pt,rounded corners=7mm]
    ([xshift=1.35cm,yshift=.85cm]current page.south west) rectangle
    ([xshift=-1.35cm,yshift=4.25cm]current page.south east);
  \node[anchor=center,text=milcNegro,text width=17.0cm,align=center] at ([yshift=2.55cm]current page.south)
    {\sffamily\fontsize{10}{12}\selectfont Metodología MILC · Apertura ancestral · Pensamiento computacional · Triángulo Dussel-Estoicismo-Floridi\\
    \textcolor{milcGradeDark}{\bfseries Producto:} 5 entrevistas a personas afectadas por el problema elegido en la sesión 1 (con guion de 7 preguntas + audio o notas + nombre y oficio del entrevistado) + síntesis de hallazgos en 1 página estructurada en 4 bloques (lo que confirmé, lo que descarté, lo que no esperaba, decisión de seguir o pivotar)};
\end{tikzpicture}
\mbox{}
\newpage

\section*{Apertura: Validación --- hablar antes de construir}

\begin{softbox}{milcGradeDark}{milcGradeSoft}{Saber ancestral}
En las veredas del Valle del Cauca y en los cabildos indígenas del Pacífico, antes de iniciar cualquier obra nueva había una costumbre obligatoria: \textbf{ir a hablar con los mayores}. Antes de sembrar un cultivo nuevo, se consultaba a la abuela campesina que conocía la luna; antes de construir una casa, se hablaba con el albañil viejo que sabía la pendiente; antes de abrir un negocio en la plaza, se preguntaba a los vendedores del lugar qué pasaba con sus clientes. No era trámite ni cortesía: era \textbf{validación de la realidad}. La persona que actuaba sin consultar a quien sabía era llamada \emph{terco} o \emph{novato}, y casi siempre fracasaba en el primer año. La sabiduría barrial y rural sostiene una verdad incómoda para el emprendedor moderno: \textbf{lo que tú crees del mundo y lo que el mundo es no coinciden hasta que preguntas}.
\end{softbox}

\begin{softbox}{milcTurquesa}{milcGris}{Saber contemporáneo}
La \textbf{validación} en metodologías contemporáneas (Customer Discovery de Steve Blank, Mom Test de Rob Fitzpatrick, Lean Startup de Eric Ries) es la práctica de \textbf{conversar con afectados reales antes de construir nada}. El error más caro del emprendedor novato es \emph{enamorarse de su hipótesis} y construir 6 meses sin haber hablado con un solo afectado. La regla profesional: \textbf{5 entrevistas valen más que 50 horas de planeación}. Una entrevista de validación bien hecha tiene 3 propiedades: (1) pregunta por \textbf{pasado concreto}, no por futuro hipotético (no \emph{``¿usarías una app que…?''} sino \emph{``cuéntame la última vez que tuviste este problema''}); (2) busca \textbf{información}, no \textbf{aprobación} (el entrevistado no debe sentir que te tiene que dar la razón); (3) deja \textbf{silencio} para que la persona se extienda (el oro está en lo que dice de más, no en lo que responde literal).
\end{softbox}

\begin{softbox}{milcMostaza}{milcGris}{Pregunta-puente}
\textit{¿Qué sabía el campesino al consultar a la abuela antes de sembrar, que el emprendedor novato olvida cuando lanza su proyecto sin hablar con un afectado? ¿Y por qué \emph{``¿usarías mi app?''} es la peor pregunta de validación posible, aunque parezca la más natural?}
\end{softbox}

\begin{softbox}{milcVerde}{milcGris}{Saber y saber hacer}
Al terminar podrás: (1) \textbf{identificar} 5 personas afectadas reales y agendarlas para entrevista esta semana; (2) \textbf{analizar} la diferencia entre preguntas malas (hipotéticas, sugerentes, de aprobación) y preguntas buenas (de pasado concreto, abiertas, neutras); (3) \textbf{explicar} qué dice el método Mom Test y cómo aplica a tu proyecto; (4) \textbf{evaluar} tus 5 entrevistas con un cuadro de hallazgos y decidir con honestidad si tu hipótesis inicial \emph{sobrevive}, \emph{pivota} o \emph{muere}.
\end{softbox}



\small
\begin{tabularx}{\textwidth}{>{\bfseries}p{3.0cm}Y}
\rowcolor{milcGradePrimary!12}Autor & Dr. Álvaro Cárdenas Orozco\\
DBA & Diseño, implementación y sustentación de un proyecto de impacto local (MEN, T\&I)\\
Referentes & Dussel (analéctica · sur global) · Lean Startup · Floridi · Estoicismo\\
Producto final & 5 entrevistas a personas afectadas por el problema elegido en la sesión 1 (con guion de 7 preguntas + audio o notas + nombre y oficio del entrevistado) + síntesis de hallazgos en 1 página estructurada en 4 bloques (lo que confirmé, lo que descarté, lo que no esperaba, decisión de seguir o pivotar)\\
\end{tabularx}
\normalsize

\puente{Antes de empezar, mira el viaje completo. En la sesión 1 elegiste un problema con ficha de oportunidad. Hoy lo \textbf{contrastas con la realidad}: 5 entrevistas a personas reales que sufren ese problema. La validación es el filtro más severo del periodo. Si tu hipótesis no sobrevive a 5 entrevistas, mejor descubrirlo hoy que en la sesión 8 cuando ya tengas prototipo.}

\begin{titlebox}{milcGradeDark}
Ruta MILC: así vamos a aprender
\end{titlebox}
\begin{tabularx}{\textwidth}{>{\centering\arraybackslash}X>{\centering\arraybackslash}X>{\centering\arraybackslash}X>{\centering\arraybackslash}X}
\phasecard{1}{milcVerde}{milcGris}{Escucha\\[-1mm]\small Observo, escucho y pregunto.} &
\phasecard{2}{milcTurquesa}{milcGris}{Sistematización\\[-1mm]\small Pensamiento computacional.} &
\phasecard{3}{milcMagenta}{milcGris}{Praxis\\[-1mm]\small Construyo y aplico.} &
\phasecard{4}{milcVino}{milcGris}{Evaluación\\[-1mm]\small Reflexión liberadora.}
\end{tabularx}


\puente{Antes de teorizar, vas a hacer la \emph{escucha del territorio}: revisar tu ficha de la sesión 1 y agendar las 5 entrevistas concretas. Vas a confirmar quiénes son, cómo los contactas y cuándo hablas con cada uno esta semana.}

\begin{titlebox}{milcVerde}
Fase 1 · Escucha
\end{titlebox}

\begin{softbox}{milcVerde}{milcGris}{Escena para escuchar}
\textbf{Actividad 1 · IDENTIFICA --- Lista cerrada de 5 entrevistados} (15 min · individual). Revisa la ficha de oportunidad de la sesión 1 y la lista de personas contactables. Cierra una lista de 5 entrevistados con: (1) nombre completo y oficio; (2) relación con el problema (\emph{lo sufre, lo ve sufrir, lo intenta resolver}); (3) canal de contacto (presencial, WhatsApp, llamada); (4) cuándo lo entrevistas esta semana (día y hora tentativa); (5) variación entre ellos: al menos 3 perfiles distintos (no entrevistar 5 amigos del mismo grupo). La diversidad fortalece la validación.
\end{softbox}

\checkbox Cerré lista de 5 entrevistados con nombre y oficio.\\
\checkbox Al menos 3 perfiles distintos (no entrevisté solo amigos).\\
\checkbox Cada entrevistado tiene día y hora tentativa esta semana.

\begin{infoband}{milcVerde}


\textbf{Lo que va al cuaderno (Actividad 1 · IDENTIFICA).} Título: «Actividad 1 --- Mis 5 entrevistados». \textbf{Formato:} tabla 5 filas × 5 columnas (Nombre + oficio~|~Relación con el problema~|~Canal~|~Día/hora~|~Por qué este perfil). \textbf{Sabes que terminaste cuando} podrías mandar los 5 mensajes de agendamiento esta misma noche.
\end{infoband}

\puente{Ya tienes los entrevistados. Ahora vas a entender la \textbf{anatomía de la entrevista profesional}: las 7 preguntas del guion, los 5 errores mortales (preguntas hipotéticas, sugerentes, de aprobación, cerradas, sobre tu idea), y la diferencia entre \emph{entrevista} y \emph{encuesta}.}

\begin{titlebox}{milcTurquesa}
Fase 2 · Sistematización con pensamiento computacional
\end{titlebox}

\begin{softbox}{milcTurquesa}{milcGris}{Cuatro pilares aplicados al tema}
Una entrevista mal hecha es peor que no entrevistar: produce \textbf{falsa validación}. Para evitarlo necesitas la anatomía profesional: 7 preguntas del Mom Test adaptadas a tu problema, y los 5 errores que matan la información antes de obtenerla.
\end{softbox}

\small
\begin{tabularx}{\textwidth}{>{\bfseries\RaggedRight\arraybackslash}p{3.6cm}Y}
\rowcolor{milcTurquesa!90}\color{white}Pilar & \color{white}\bfseries Aplicación al tema\\
1. Descomposición & El guion profesional se descompone en 7 preguntas: (1) \textbf{Contexto}: \emph{``cuéntame de tu día/semana en relación con [tema del problema]''}. (2) \textbf{Última vez}: \emph{``¿cuándo fue la última vez que te pasó X?, cuéntame cómo fue''}. (3) \textbf{Qué hiciste}: \emph{``¿qué hiciste tú en ese momento?''} (revela el apaño actual). (4) \textbf{Qué intentaste antes}: \emph{``¿qué soluciones probaste y por qué no funcionaron?''}. (5) \textbf{Cuánto te costó}: \emph{``¿cuánto tiempo o dinero te costó esa última vez?''} (revela el tamaño del dolor). (6) \textbf{Qué te frenó}: \emph{``¿qué te impidió resolverlo bien?''}. (7) \textbf{Quién más}: \emph{``¿conoces a alguien más que viva lo mismo? ¿podrías presentármelo?''} (cadena de validación).\\
2. Reconocimiento de patrones & Los 5 errores mortales del entrevistador novato son: (a) \textbf{Pregunta hipotética}: \emph{``¿usarías una app que…?''} (las personas mienten sobre el futuro, no sobre el pasado). (b) \textbf{Pregunta sugerente}: \emph{``¿no te parece terrible que…?''} (induce la respuesta). (c) \textbf{Pregunta sobre tu idea}: \emph{``¿qué te parece mi solución?''} (la persona no tiene contexto y dirá algo amable). (d) \textbf{Pregunta cerrada}: \emph{``¿estás de acuerdo con…?''} (cierra el flujo en sí o no). (e) \textbf{Pregunta colectiva}: \emph{``¿qué hace la gente cuando…?''} (los individuos no saben qué hace \emph{la gente}; saben qué hicieron ellos).\\
3. Abstracción & Lo esencial cabe en una frase: \textbf{pregunta por hechos del pasado, no por opiniones del futuro}. Steve Blank lo llamó \emph{``get out of the building''}; Rob Fitzpatrick lo llamó \emph{``Mom Test''} (preguntas tan bien hechas que ni siquiera tu mamá te puede mentir respondiéndolas). Es el mismo principio: la verdad vive en el pasado concreto, la ilusión en el futuro hipotético.\\
4. Algoritmo conceptual & Algoritmo: (1) adapta las 7 preguntas a tu problema específico; (2) léelas en voz alta y elimina las que sean hipotéticas o sugerentes; (3) prepara una matriz de hallazgos vacía con 7 columnas (una por pregunta) y 5 filas (una por entrevistado); (4) durante cada entrevista llena la matriz en tiempo real; (5) después de cada entrevista, anota 1-2 frases textuales del entrevistado que te sorprendieron.\\
\end{tabularx}
\normalsize

\begin{drawbox}{milcTurquesa}{Anatomía de la entrevista de validación}
\textbf{Contexto:} día o semana de la persona en relación al tema. \\[1mm] \textbf{Última vez:} narrativa concreta del último episodio. \\[1mm] \textbf{Qué hiciste:} apaño actual real. \\[1mm] \textbf{Qué intentaste antes:} soluciones fallidas. \\[1mm] \textbf{Cuánto te costó:} tiempo o dinero del último episodio. \\[1mm] \textbf{Qué te frenó:} barrera real (no imaginada). \\[1mm] \textbf{Quién más:} cadena de referidos.
\end{drawbox}

\begin{softbox}{milcMostaza}{milcGris}{Errores comunes y su corrección}
(1) Preguntar \emph{``¿usarías…?''} en lugar de \emph{``¿cuándo fue la última vez…?''}. (2) Hablar más que el entrevistado (regla profesional: el entrevistado habla 80\%, tú 20\%). (3) Defender tu idea cuando aparece evidencia contraria. (4) Quedarse en preguntas demográficas (edad, ciudad) sin entrar al dolor. (5) No pedir referidos al final (perder la cadena de validación). (6) Entrevistar solo a amigos que te dirán que sí (sesgo de complacencia). (7) Anotar después de la entrevista en lugar de durante (se pierde el 70\% del detalle).
\end{softbox}

\begin{infoband}{milcTurquesa}


\textbf{Lo que va al cuaderno (Actividad 2 · ANALIZA + EXPLICA).} Título: «Actividad 2 --- Anatomía de la entrevista». \textbf{Formato:} (a) las 7 preguntas adaptadas a tu problema; (b) los 5 errores mortales con marca de cuál te tienta más; (c) matriz de hallazgos vacía (5 filas × 7 columnas). \textbf{Sabes que terminaste cuando} tienes el guion impreso o en celular, listo para usar en la primera entrevista.
\end{infoband}

\puente{Tienes el método. Toca aplicarlo: redactar el guion de 7 preguntas adaptado a tu problema y producir la matriz de hallazgos donde anotarás lo dicho por cada entrevistado. Las entrevistas se hacen durante la semana; la sesión termina con el material listo.}

\begin{titlebox}{milcMagenta}
Fase 3 · Praxis con pensamiento computacional
\end{titlebox}

\begin{softbox}{milcMagenta}{milcGris}{Construyendo el producto con los cuatro pilares}
\textbf{Actividad 3 · EVALÚA --- Síntesis de hallazgos y decisión} (30 min · individual, después de las 5 entrevistas). Esta actividad se completa al final de la semana, cuando ya entrevistaste a las 5 personas. Aquí no se inventa: se sintetiza la matriz, se nombran los hallazgos y se toma una decisión con honestidad: \emph{sigo, pivoto o descarto}.
\end{softbox}

\small
\begin{tabularx}{\textwidth}{>{\bfseries\RaggedRight\arraybackslash}p{3.6cm}Y}
\rowcolor{milcMagenta!90}\color{white}Pilar & \color{white}\bfseries Aplicación al producto\\
Descomposición & Tu síntesis se estructura en 4 bloques: (1) \textbf{Lo que confirmé}: hipótesis de la sesión 1 que sobrevivieron a las 5 entrevistas (ejemplo: \emph{``confirmo que los tenderos sí pierden ventas por no avisar a clientes regulares''}). (2) \textbf{Lo que descarté}: hipótesis que el campo rechazó (ejemplo: \emph{``creía que el problema era el inventario, pero ningún entrevistado lo mencionó: era la comunicación''}). (3) \textbf{Lo que no esperaba}: hallazgos nuevos no previstos en la ficha original (las mejores entrevistas siempre tienen sorpresas). (4) \textbf{Decisión}: \emph{sigo} (la hipótesis sobrevivió), \emph{pivoto} (cambio el problema o el afectado), \emph{descarto} (el problema no existe como yo lo describí).\\
Patrones & Sigue el patrón \textbf{``cita textual sobre interpretación''}: en la síntesis, prioriza frases textuales de los entrevistados sobre tus interpretaciones. Si 3 de 5 dijeron \emph{``ya intenté hacer eso y no funcionó''}, esa frase vale más que tu párrafo explicando por qué crees que sí funcionaría. La phronesis del emprendedor consiste en \textbf{escuchar lo que el dato dice, no lo que uno quiere que diga}.\\
Abstracción & Pivotar no es fracasar: es la decisión profesional más respetada en el método Lean. Pivotar a tiempo (sesión 2) ahorra 8 sesiones de trabajo construido sobre supuestos falsos. La phronesis del oficio respeta a quien pivota con evidencia más que a quien sostiene su idea original por orgullo. Tu rúbrica del periodo valora más una decisión de pivote bien argumentada que una hipótesis original mal validada.\\
Algoritmo de armado & Algoritmo: (1) recopila las 5 matrices llenas; (2) marca con highlighter las frases textuales más fuertes; (3) cuenta cuántos entrevistados mencionaron espontáneamente cada dolor (la frecuencia es el dato más importante); (4) redacta los 4 bloques de síntesis en 1 página; (5) toma la decisión por escrito con justificación de 2 líneas; (6) si pivotas, redacta la nueva ficha de oportunidad antes de cerrar la sesión.\\
\end{tabularx}
\normalsize

\begin{drawbox}{milcMagenta}{Checklist antes de entregar la síntesis}
\checkbox 5 entrevistas realizadas a 5 personas distintas. \\[1mm] \checkbox Matriz de hallazgos llena durante cada entrevista (no después). \\[1mm] \checkbox Citas textuales destacadas de cada entrevistado. \\[1mm] \checkbox Síntesis en 4 bloques (confirmé, descarté, no esperaba, decisión). \\[1mm] \checkbox Decisión explícita: sigo, pivoto o descarto. \\[1mm] \checkbox Si pivotas, ficha de oportunidad actualizada. \\[1mm] \checkbox Síntesis cabe en 1 página de cuaderno o documento.
\end{drawbox}

\begin{softbox}{milcMagenta}{milcGris}{Plantilla de trabajo}
\textbf{Lo que confirmé.} \textit{Hipótesis de S1 que sobrevivieron + cuántos entrevistados las apoyaron.} \\[1mm] \textbf{Lo que descarté.} \textit{Hipótesis que el campo rechazó + evidencia textual.} \\[1mm] \textbf{Lo que no esperaba.} \textit{Hallazgos nuevos + 1 cita textual por hallazgo.} \\[1mm] \textbf{Decisión.} \textit{Sigo / Pivoto / Descarto + justificación de 2 líneas.}
\end{softbox}

\begin{infoband}{milcMagenta}


\textbf{Lo que va al cuaderno (Actividad 3 · EVALÚA).} Título: «Actividad 3 --- Síntesis de validación». \textbf{Formato:} 1 página con 4 bloques de síntesis + decisión + (si pivotas) nueva ficha de oportunidad. \textbf{Sabes que terminaste cuando} podrías presentar la decisión a un asesor externo y defenderla con citas textuales en mano.
\end{infoband}

\puente{Lo que sigue resume \emph{qué} entregas al final de la semana: 5 entrevistas reales documentadas + síntesis de hallazgos en 1 página + decisión de seguir, pivotar o descartar el problema. El criterio principal: que un consultor leyendo tu síntesis pueda decir \emph{``con esta evidencia, la decisión está bien tomada''}, sin tener que confiar en tu palabra.}

\begin{titlebox}{milcMostaza}
Producto de la sesión
\end{titlebox}
\textbf{5 entrevistas + matriz de hallazgos + síntesis con decisión}.
\begin{softbox}{milcMostaza}{milcGris}{Criterios de calidad}
(1) 5 entrevistas a personas reales y diversas. (2) Guion de 7 preguntas sin errores mortales. (3) Matriz de hallazgos llena durante cada entrevista. (4) Síntesis en 4 bloques con citas textuales. (5) Decisión explícita (sigo, pivoto, descarto) con justificación. (6) Si pivota, ficha de oportunidad actualizada.
\end{softbox}

\puente{Antes de cerrar, mira la validación desde las cinco dimensiones humanas. Conversar con afectados reales es respeto por el oficio; saltarse las entrevistas y inventarlas es la manera más rápida de perder un periodo entero.}

\begin{titlebox}{milcVino}
Fase 4 · Evaluación liberadora — cinco dimensiones
\end{titlebox}

\begin{softbox}{milcVino}{milcGris}{Sentido de la evaluación}
Evaluar no es solo revisar una nota. Es reconocer qué aprendí, qué cambió en mí, qué emociones manejé, cómo me relaciono con la ciudad y la comunidad, y qué le dejaré a quien viene detrás.
\end{softbox}

\textbf{Mi reflexión liberadora en cinco dimensiones.} Completa cada frase iniciada:

\small
\begin{tabularx}{\textwidth}{>{\bfseries\RaggedRight\arraybackslash}p{3.4cm}Y>{\RaggedRight\arraybackslash}p{4.6cm}}
\rowcolor{milcVino}\color{white}Dimensión & \color{white}\bfseries Mi respuesta (frame starter) & \color{white}\bfseries Algo que descubrí\\
1. Personal & Lo que más cambió en mí fue \linead{6.5cm} & \linead{4.2cm}\\
2. Emocional & La emoción más fuerte fue \linead{2.5cm} y la regulé \linead{3cm} & \linead{4.2cm}\\
3. Ciudadana & En democracia esto importa porque \linead{6cm} & \linead{4.2cm}\\
4. Local (Cartago) & En mi barrio sería útil para \linead{6.5cm} & \linead{4.2cm}\\
5. Intergeneracional & A mi abuela/abuelo le diría \linead{6.5cm} & \linead{4.2cm}\\
\end{tabularx}
\normalsize

\begin{infoband}{milcVino}
\textbf{Para la próxima sustentación.} Vuelve a esta tabla en seis meses. Verás que las respuestas cambian — no porque lo aprendido cambie, sino porque tú cambias. La evaluación liberadora es una conversación contigo a través del tiempo.
\end{infoband}


\puente{Y para terminar, tres voces sobre la escucha y la verdad. Dussel sobre el lugar epistémico del afectado, Séneca sobre la valentía de aceptar evidencia contraria, Floridi sobre la ética de tomar decisiones sin datos cuando los datos están al alcance.}

\begin{titlebox}{milcGradeDark}
Triángulo de pensamiento · cierre filosófico
\end{titlebox}

\begin{softbox}{milcVino}{milcGris}{Dussel · lente del nosotros}
\textit{``El afectado es el primer epistemólogo del problema; nadie sabe del dolor mejor que quien lo sufre.''}

Tu validación es ética antes de ser metodológica. El emprendedor que decide \emph{por} el afectado, sin escucharlo primero, repite la lógica colonial de quien sabe más que el pueblo. La filosofía de la liberación pide lo opuesto: \textbf{el lugar epistémico privilegiado es el del afectado}, no el del que diseña la solución. Cuando entrevistas a tus 5 personas con respeto y silencio, no estás haciendo trámite metodológico: estás reconociendo que ellas saben de su problema más que tú.

\textbf{En mis entrevistas, ¿escuché de verdad o solo busqué confirmación? ¿Hablé yo más o habló el entrevistado? ¿Acepté lo que dijo cuando contradijo mi idea?}
\end{softbox}
\linead{\linewidth}\\[1mm]
\linead{\linewidth}

\begin{softbox}{milcMostaza}{milcGris}{Estoicismo · Séneca · cuidado interior}
\textit{``La evidencia contraria es regalo, no insulto. El sabio cambia de opinión cuando los datos lo exigen.''}

Es tentador defender tu hipótesis cuando un entrevistado la contradice. La sabiduría estoica recomienda lo opuesto: \emph{recibe la evidencia contraria como regalo}. Cambiar de opinión por datos no es debilidad: es disciplina racional. El emprendedor que sostiene su idea contra 5 entrevistas adversas no muestra valor: muestra terquedad. La phronesis del oficio honra a quien pivota con evidencia más que a quien insiste sin ella.

\textbf{Si mis 5 entrevistas contradijeran mi hipótesis, ¿tendría la valentía de pivotar hoy, o seguiría adelante por orgullo?}
\end{softbox}
\linead{\linewidth}\\[1mm]
\linead{\linewidth}

\begin{softbox}{milcTurquesa}{milcGris}{Floridi · lente de la infoesfera}
\textit{``Decidir sin datos cuando los datos están al alcance es la nueva forma de irresponsabilidad profesional.''}

Antes de internet, decidir sin datos era inevitable: la información era costosa. Hoy es lo opuesto: hablar con 5 personas reales cuesta una semana, y la mayoría de los emprendedores no lo hace. La ética digital contemporánea declara esa omisión como irresponsable. El profesional que invierte horas en una solución sin haber hablado con un afectado tira a la basura el principio básico del oficio: \textbf{los datos primero, las decisiones después}.

\textbf{¿Estoy decidiendo mi proyecto con base en 5 conversaciones reales, o con base en lo que yo imagino del mundo?}
\end{softbox}
\linead{\linewidth}\\[1mm]
\linead{\linewidth}

\begin{titlebox}{milcVino}
Compromiso final
\end{titlebox}

\small
\begin{tabularx}{\textwidth}{>{\RaggedRight\arraybackslash}p{3.0cm}YYY}
\rowcolor{milcVino}\color{white}\bfseries Criterio & \color{white}\bfseries Lo logré & \color{white}\bfseries En proceso & \color{white}\bfseries Necesito apoyo\\
Comprensión del tema & Explico ideas centrales con mis palabras. & Reconozco ideas pero falta precisión. & Necesito repasar conceptos base.\\
Pensamiento computacional & Apliqué los 4 pilares al diseñar mi producto. & Apliqué algunos pilares. & Necesito releer la fase 2.\\
Aplicación y producto & Mi producto cumple los criterios y se puede revisar. & Producto funcional con ajustes pendientes. & Aún en construcción.\\
Reflexión 5 dimensiones & Respondí cada dimensión con honestidad. & Respondí algunas con detalle. & Mis respuestas son muy generales.\\
Triángulo de pensamiento & Conecté las 3 voces con mi trabajo real. & Conecté al menos una. & No relaciono las voces con mi proceso.\\
\end{tabularx}
\normalsize

\textbf{Hoy entendí que:}\\[1mm]
\linead{\linewidth}\\[1mm]
\linead{\linewidth}

\textbf{Mi compromiso para la próxima sesión es:}\\[1mm]
\linead{\linewidth}\\[1mm]
\linead{\linewidth}

\begin{softbox}{milcMostaza}{milcGris}{Cita de despedida · Marco Aurelio}
\textit{``El obstáculo en el camino se vuelve el camino. Lo que estorba al camino se vuelve camino.''}\\[1mm]
Cada error de hoy, bien observado, se convierte en pieza del camino que pisarás mañana.
\end{softbox}

\small
\begin{tabularx}{\textwidth}{>{\bfseries}p{3.6cm}Y}
\rowcolor{milcGradeDark!12}Nombre completo: & \linead{10cm}\\
Fecha: & \linead{4cm} \quad Curso/grupo: \linead{4cm}\\
Mi firma: & \linead{9cm}\\
\end{tabularx}
\normalsize

\begin{infoband}{milcGradeDark}
\textbf{Para la próxima sesión.} Llega con las 5 entrevistas hechas y la síntesis lista. En la sesión 3 vas a transformar lo validado en \textbf{modelo de negocio} con el Business Model Canvas. Si pivotaste, llega con la ficha de oportunidad actualizada. La validación se honra construyendo sobre ella, no archivándola.
\end{infoband}

\end{document}
